%%% Докторская диссертация — полный шаблон (все элементы)
%%% Тематика: Стохастическое управление
%%% Компиляция: xelatex → biber → xelatex → xelatex

\documentclass[font=modern]{cmcmsuthesis-doctoral}

\author{Петров Сергей Николаевич}
\title{Теория устойчивости стохастических систем
       с марковскими переключениями}
\thesissetup{department={Кафедра теории вероятностей}}
\supervisor{name={академик РАН, д-р~физ.-мат.~наук, проф. Иванов~И.И.}}
\specialty{code={1.1.5},name={математическая физика}}

\usepackage{biblatex}
\addbibresource{refs.bib}

\begin{document}
\maketitle
\tableofcontents

% ================================================================
\chapter*{Введение}
\addcontentsline{toc}{chapter}{Введение}

Стохастические системы с марковскими переключениями
возникают в управлении~\cite{shannon1948,macwilliams1979}.

\textbf{Целью} работы является развитие теории устойчивости.

\textbf{Задачи}:
\begin{enumerate}[beginpenalty=10000]
    \item получить критерий экспоненциальной устойчивости;
    \item построить методы синтеза управлений;
    \item обобщить на случай запаздывания;
    \item исследовать робастную устойчивость.
\end{enumerate}

% ================================================================
\chapter{Устойчивость линейных систем}\label{ch:stability}

\section{Основные определения}

Пусть $\vect{x}(t)\in\mathbb{R}^n$, $\xi_t$~--- марковская цепь.
$\matr{A}\in\mathbb{F}^{k\times n}_q$, $\str{S}_n$,
$\elt{\sigma}$, $\code{C}$, $\algo{LMI}$,
$\compl{O}(n^{3.5})$, $\crypto{E}$.

\begin{definition}\label{def:system}
\emph{Стохастической системой с переключениями} называется
\begin{equation}\label{eq:sde}
d\vect{x}=\matr{A}(\xi_t)\vect{x}\,dt+\matr{B}(\xi_t)\vect{x}\,dW.
\end{equation}
\end{definition}

\begin{definition}\label{def:ms-stable}
Система~\eqref{eq:sde} \emph{устойчива в среднеквадратичном},
если $\lim_{t\to\infty}\mathbb{E}\|\vect{x}(t)\|^2=0$.
\end{definition}

\begin{theorem}\label{thm:lmi}
Система~\eqref{eq:sde} экспоненциально устойчива $\Leftrightarrow$
$\exists\matr{P}_i\succ 0$:
\[\matr{A}_i^\top\matr{P}_i+\matr{P}_i\matr{A}_i
+\matr{B}_i^\top\matr{P}_i\matr{B}_i+\sum_j q_{ij}\matr{P}_j\prec 0.\]
\end{theorem}

\begin{proof}
Функция Ляпунова $V(\vect{x},i)=\vect{x}^\top\matr{P}_i\vect{x}$
и формула Ито.
\end{proof}

\begin{corollary}\label{cor:sdp}
LMI проверяется за $\compl{O}(N\cdot n^{3.5})$.
\end{corollary}

\begin{lemma}\label{lem:lyapunov}
Если $\mathcal{L}V\leq-\alpha V$, то
$\mathbb{E}V(t)\leq e^{-\alpha t}V(0)$.
\end{lemma}

\begin{proposition}\label{prop:necessary}
Если система устойчива, то $\operatorname{Re}\lambda_i(\matr{A}_j)<0$
для некоторого $j$.
\end{proposition}

\begin{remark}\label{rem:switching}
Устойчивость каждой подсистемы не гарантирует устойчивости
системы с переключениями.
\end{remark}

\begin{example}\label{ex:two-mode}
Двухрежимная система: $\matr{A}_1=\begin{pmatrix}-1&0\\0&-2\end{pmatrix}$,
$\matr{A}_2=\begin{pmatrix}-0.5&1\\-1&-0.5\end{pmatrix}$.
\end{example}

\begin{problem}\label{prob:delay}
Обобщить критерий теоремы~\ref{thm:lmi} на системы с запаздыванием.
\end{problem}

% ================================================================
\chapter{Формулы}\label{ch:formulas}

\section{Ненумерованные формулы}
Inline: $x\approx\sin x$ при $x\to 0$.
Display: \[(x_1+x_2)^2=x_1^2+2x_1x_2+x_2^2.\]
Непрерывная дробь:
\[\frac{1}{\sqrt{2}+\displaystyle\frac{1}{\sqrt{2}+\displaystyle\frac{1}{\sqrt{2}+\cdots}}}\]

\section{Нумерованные формулы}
\begin{equation}\label{eq:euler}e=\lim_{n\to\infty}\left(1+\frac{1}{n}\right)^{\!n}.\end{equation}
\begin{equation}\label{eq:basel}\lim_{n\to\infty}\sum_{k=1}^n\frac{1}{k^2}=\frac{\pi^2}{6}.\end{equation}
Ссылки: \eqref{eq:euler} и~\eqref{eq:basel}.
\begin{equation}\label{eq:long}\begin{multlined}1+2+3+\dots+50+\dots+\\+96+97+98+99+100=5050.\end{multlined}\end{equation}

\section{Многострочные формулы}
\[\begin{aligned}f_W&=\min\!\left(1,\max\!\left(0,\frac{W/W_{\max}}{W_{\text{crit}}}\right)\right),\\f_T&=\min\!\left(1,\max\!\left(0,\frac{T_s/T_{\text{melt}}}{T_{\text{crit}}}\right)\right).\end{aligned}\]

\[|x|=\left\{\begin{alignedat}{2}& &x,\quad&\text{если }x\geqslant 0;\\&-&x,\quad&\text{если }x<0.\end{alignedat}\right.\]

\[|x|=\begin{cases}\phantom{-}x,&\text{если }x\geqslant 0;\\-x,&\text{если }x<0.\end{cases}\]

\begin{equation}\label{eq:ito}\begin{alignedat}{2}
d\mathbb{E}V&=-\alpha\,\mathbb{E}V\,dt,&\quad&\\
\mathbb{E}V(t)&=e^{-\alpha t}\,V(0).&&
\end{alignedat}\end{equation}

\section{Субформулы}
\begin{subequations}\label{eq:sub1}\begin{gather}y=x^2+1\label{eq:s1a}\\y=2x^2-x+1\label{eq:s1b}\\y=3x^2-4x+8\label{eq:s1c}\end{gather}\end{subequations}
\begin{subequations}\label{eq:sub2}\begin{align}y&=x^3+x^2+x+1\label{eq:s2a}\\y&=x^2\label{eq:s2b}\end{align}\end{subequations}
Ссылки: \eqref{eq:s1a}, \eqref{eq:s1b}, \eqref{eq:s2a}, \eqref{eq:s1c}.

Система (Лоренц): \[\left\{\begin{array}{rl}\dot x=&\sigma(y-x)\\\dot y=&x(r-z)-y\\\dot z=&xy-bz\end{array}\right..\]

Матрица: \[\left(\begin{array}{ccc}a_{11}&\ldots&a_{1n}\\\vdots&\ddots&\vdots\\a_{n1}&\ldots&a_{nn}\end{array}\right).\]

Длинная inline: $a_1+a_2+a_3+a_4+a_5+a_6+a_7+a_8+a_9+a_{10}+a_{11}+a_{12}+a_{13}+a_{14}+a_{15}+a_{16}+a_{17}+a_{18}+a_{19}+a_{20}=\sum_{i=1}^{20}a_i.$

\section{Греческие буквы и алфавиты}
Обычные: $\alpha\beta\gamma\delta\epsilon\theta\lambda\mu\pi\sigma\phi\omega\Gamma\Delta\Omega$.
Прямые: $\symup{\alpha}\symup{\beta}\symup{\gamma}\symup{\delta}\symup{\Gamma}\symup{\Omega}$.
Жирные: $\symbf{\alpha}\symbf{\beta}\symbf{\gamma}\symbf{\Gamma}\symbf{\Omega}$.
Жирные прямые: $\symbfup{\alpha}\symbfup{\beta}\symbfup{\gamma}\symbfup{\Gamma}\symbfup{\Omega}$.

\[\begin{aligned}\mathcal{ABCDEFGHIJKLMNOPQRSTUVWXYZ}\\\mathfrak{ABCDEFGHIJKLMNOPQRSTUVWXYZ}\\\mathbb{ABCDEFGHIJKLMNOPQRSTUVWXYZ}\end{aligned}\]

Русские операторы: $\tan\alpha$/$\cot\beta$/$\csc\gamma$; $\le$/$\leqslant$; $\phi$/$\varphi$; $\epsilon$/$\varepsilon$; $\kappa$/$\varkappa$; $\emptyset$/$\varnothing$.

Отступы:
\[\begin{aligned}\backslash!&\quad f(x)=x^2\!+3x\!+2\\\text{default}&\quad f(x)=x^2+3x+2\\\backslash,&\quad f(x)=x^2\,+3x\,+2\\\backslash;&\quad f(x)=x^2\;+3x\;+2\\\backslash\text{quad}&\quad f(x)=x^2\quad+3x\quad+2\end{aligned}\]

% ================================================================
\chapter{Форматирование, таблицы, изображения и списки}

\section{Форматирование текста}
\begin{description}
\item[Жирный:] \textbf{жирный}.\item[Курсив:] \emph{курсив}.
\item[Жирный курсив:] \textbf{\emph{жирный курсив}}.
\item[Подчёркивание:] \underline{подчёркнутый}.
\item[Зачёркивание:] \sout{зачёркнутый}.
\item[Sans:] \textsf{sans serif font}.
\end{description}

\section{Таблицы}
\begin{table}[htbp]\caption{Простая таблица}\label{tab:simple}\centering
\begin{tabular}{|l|c|c|c|}\hline Система&$n$&$N$&Устойчива\\\hline
Двухрежимная&2&2&да\\\hline Трёхрежимная&3&3&нет\\\hline
С запаздыванием&4&2&да\\\hline
\end{tabular}\end{table}

\begin{table}[htbp]\caption{Booktabs}\label{tab:booktabs}\centering
\begin{tabular}{lrcc}\toprule
\textbf{$n$}&\textbf{$N$}&\textbf{Время SDP,с}&\textbf{$\lambda_{\max}$}\\\midrule
2&3&0.01&$-0.42$\\5&4&0.08&$-0.31$\\10&5&0.95&$-0.18$\\20&8&12.4&$-0.09$\\\bottomrule
\end{tabular}\end{table}

\begin{longtable}{@{}rr@{}}\caption{Длинная таблица}\label{tab:long}\\
\toprule $t$&$\mathbb{E}\|\vect{x}\|^2$\\\midrule\endfirsthead
\toprule $t$&$\mathbb{E}\|\vect{x}\|^2$\\\midrule\endhead\bottomrule\endfoot
0.0&1.000\\0.5&0.657\\1.0&0.432\\1.5&0.284\\2.0&0.187\\
2.5&0.123\\3.0&0.081\\3.5&0.053\\4.0&0.035\\4.5&0.023\\\end{longtable}
Ссылка: таблица~\ref{tab:booktabs}.

\section{Изображения}
%% \begin{figure}[htbp]\centering\includegraphics[width=0.8\textwidth]{images/trajectories.png}\caption{Траектории}\label{fig:traj}\end{figure}
%% \begin{figure}[htbp]\centering\includegraphics[scale=0.5]{images/lyapunov.png}\caption{Функция Ляпунова (масштаб 0.5)}\label{fig:lyap}\end{figure}
%% Ссылка: рисунок~\ref{fig:traj}.

\section{Списки}
\subsection{Нумерованный}
\begin{enumerate}\item Первый\begin{enumerate}\item Второй\begin{enumerate}\item Третий\begin{enumerate}\item Четвёртый\end{enumerate}\end{enumerate}\end{enumerate}\end{enumerate}
\subsection{Ненумерованный}
\begin{itemize}\item Первый\begin{itemize}\item Второй\begin{itemize}\item Третий\begin{itemize}\item Четвёртый\end{itemize}\end{itemize}\end{itemize}\end{itemize}

% ================================================================
\chapter{Алгоритмы и листинги}

\begin{algorithm}[H]\caption{Синтез стабилизации}\label{alg:synth}
\KwInput{$\matr{A}_i,\matr{B}_i,\matr{D}_i$, генератор $\matr{Q}$}
\KwOutput{$\matr{K}_i$ или «невозможно»}
Сформировать LMI\;
Решить SDP\tcp*{полуопределённое программирование}
\Comment{проверяем решение}
\uIf{разрешимо}{
    \For{$i=1$ \KwTo $N$}{$\matr{K}_i\gets\matr{Y}_i\matr{X}_i^{-1}$\;}
    \Return{$\matr{K}_1,\ldots,\matr{K}_N$}\;
}
\uElseIf{граница достигнута}{\Return{«на границе»}\;}
\Else{\Return{«невозможно»}\;}
\end{algorithm}
Ссылка: алгоритм~\ref{alg:synth}.

\begin{lstlisting}[language=Python,caption={Python},label={lst:py}]
import numpy as np
from scipy.linalg import solve_lyapunov
def check_stability(A, Q_gen):
    """Проверка устойчивости через Ляпунова."""
    n = A[0].shape[0]
    # Решаем уравнение Ляпунова
    P = solve_lyapunov(A[0].T, -np.eye(n))
    stable = np.all(np.linalg.eigvals(P) > 0)
    return stable, P
A = [np.array([[-1, 0.5], [-0.5, -2]])]
ok, P = check_stability(A, None)
print("Устойчива:", ok)  # Ожидается True
\end{lstlisting}

\begin{lstlisting}[language=C,caption={C},label={lst:c}]
#include <stdio.h>
#include <math.h>
/* Моделирование траектории */
void simulate(double *x, double A[2][2], double dt, int steps) {
    // Метод Эйлера--Маруямы
    for (int i = 0; i < steps; ++i) {
        double dx0 = A[0][0]*x[0] + A[0][1]*x[1];
        double dx1 = A[1][0]*x[0] + A[1][1]*x[1];
        x[0] += dx0*dt; x[1] += dx1*dt;
    }
}
int main(void) {
    double x[2] = {1.0, 0.5};
    double A[2][2] = {{-1,0},{0,-2}};
    simulate(x, A, 0.01, 1000);
    printf("x = (%.4f, %.4f)\n", x[0], x[1]);
    return 0;
}
\end{lstlisting}

\begin{lstlisting}[style=pseudocode,caption={Псевдокод}]
(*@\lstInput@*): матрицы $\matr{A}_i$, генератор $\matr{Q}$
(*@\lstOutput@*): устойчива / неустойчива
procedure CheckStability($\{\matr{A}_i\}$, $\matr{Q}$)
begin
    Построить LMI из теоремы
    if SDP разрешимо then
        return "устойчива"
    else
        return "неустойчива"
end
\end{lstlisting}

\begin{lstpseudocode}[caption={Монте-Карло},label={alg:mc}]
(*@\lstInput@*): система, число траекторий $M$, горизонт $T$
(*@\lstOutput@*): оценка $\mathbb{E}\|\vect{x}(T)\|^2$
function MonteCarlo($M$, $T$)
begin
    $S \gets 0$
    for $j = 1$ to $M$ do
        Смоделировать траекторию $\vect{x}^{(j)}(T)$
        $S \gets S + \|\vect{x}^{(j)}(T)\|^2$
    end
    return $S / M$
end
\end{lstpseudocode}

% ================================================================
\chapter{Библиография и ссылки}
Одиночная~\cite{shannon1948}. Книга~\cite{knuth1984}.
Множественная~\cite{shannon1948,knuth1984,macwilliams1979}.
Ссылки: теорема~\ref{thm:lmi}, определение~\ref{def:system},
формула~\eqref{eq:sde}, таблица~\ref{tab:booktabs},
алгоритм~\ref{alg:synth}.

% ================================================================
\chapter*{Заключение}
\addcontentsline{toc}{chapter}{Заключение}
\begin{enumerate}
\item Получен критерий устойчивости (теорема~\ref{thm:lmi}).
\item Разработан метод синтеза (алгоритм~\ref{alg:synth}).
\item Проведены эксперименты (таблица~\ref{tab:booktabs}).
\end{enumerate}

\clearpage\listoftables\addcontentsline{toc}{chapter}{Список таблиц}
\printbibliography[heading=bibintoc]

\appendix
\chapter{Доказательства лемм}
\section{Лемма Ляпунова}
Доказательство леммы~\ref{lem:lyapunov}.
\section{Оценки скорости убывания}
Дополнительные выкладки.

\chapter{Параметры экспериментов}
Полные данные.
\section{Двухрежимные системы}
Таблицы.

\end{document}
